\section{Inner Meaning of Augustinianism}

Before we proceed along the Arcana of the \emph{Meditations}\footnote{\url{http://www.gornahoor.net/library/Meditations\_versionTwo.pdf}}, it might be helpful to point out a link between ``gnosis'' and the ``arcana''. In the very first chapter of his \emph{Meditations} Tomberg highlights the truth that arcana represent ``enzymes'' or ``ferment'' which make the soul fruitful by awakening \emph{the deeper, ever deeper layers}, and which lead to a spiritual ``event'', which is then termed the ``Mystery''. Human secrets are merely parts of those two processes which are (for good or bad reasons) deliberately concealed, so just as arcana is higher than secret (being the ``life'' of the ``secret'') so are mysteries greater than arcana. If you are following the Gnosis seminar or read carefully, you can compare the ``ferment'' of the enzymes with the ``friction'' which Mouravieff cites as producing ``heat and light'' which can unify (and fructify) the soul.

They are speaking about the same thing, but from slightly different perspectives. If God is the circle whose center is everywhere, and circumference is nowhere (and this was the Medieval tradition in accord with Tradition), and if each of us ``fell'' in a different way and direction, then we all work our way back to God from a different point and in a different direction. Likewise, in esoteric studies, slightly different starting points and perspectives produce a seeming difference, which is on closer inspection the same thing in a different garb or mood or tone: initiates recognize each other, because they have all returned (are returning) to the center.

Mouravieff teaches that each of us starts as a conglomerate of ``little i s'', an abyss or void, analogous to the waters in Genesis. The real self has withdrawn to create a void, and then must re-emerge to re-create the chaos and triumph over the chaos. It is in this way that matter is drawn up into Spirit. These ``little i s'' reflect the random A influences which operate to keep them under General Law (Accident). One personality wants one thing, another personality wants another, and each one takes turn pushing the human entity in opposing directions, according to the randomness that operates within this magic circle (as it were), where the different directions cancel one another out. \emph{This explains why things are ``real'' but not real at the same time}: each influence is powerful in its immediacy and its proper vicinity, but loses power as it pushes the Ego out of the vector of its power, which is inversely proportional to the distance. The Ego then comes under another vector, pushing in (likely) an opposing direction. Most men spend their lives bouncing around like a pinball, and at the end of their lives, \emph{the meaning of that existence sums to an integral that approaches a Zero}. Tomberg addresses this by pointing out that the unforgivable sin (or at least the most serious) is spiritual self-complacency: most men simply accept this situation as normal. We live in a civilization where this is largely accepted as the best, maybe the only, way of conducting one's life.

It is unfashionable (for instance) to observe Lent, to read antique European writers meditatively, to honor your father and mother, to remain monogamous and faithful to your wedding vows, etc. These practices were marks of an earlier situation in which men had agreed that the randomness of Life would be subjected to a higher Law that would begin to transform the otherwise iron-clad law of general Accident. Today, the absurdity of Life is the subject of cutting-edge drama which either justifies it, celebrates it, or wallows in it: Western man no longer believes in a method of transcending this human condition. They simply can't be bothered.

Just as Mouravieff teaches the idea that ``one should get comfortable with being uncomfortable'' (credit to a participant in the seminar) in order to sustain the necessary and painful efforts to begin to change this situation, Tomberg teaches that the magician must turn ``work into play'': this is the first condition, the primal Arcana, the \emph{unum necessarium}. Ousepensky remarks that we find motives all the time to overcome inertia for worldly affairs, and so that it is of course possible (even more possible) to find motives to sustain such an effort. In older civilizations, it was the desire not to be damned, or the aspiration to achieve a ``high culture'', or certain interests that centered around clues discerned in mundane life that hinted at something ``more'' (this would be an example of discerning B influences correctly, the influences from Higher Mind that are cast into the void of A influences, and which all flow in the same direction).

And why not? Where is it written in the conscience that man should pursue complacency?

I am reading two books on Saint Augustine, one by F Seraphim Rose\footnote{\url{https://www.amazon.com/Blessed-Augustine-Orthodox-Church-Theological/dp/0938635123/ref=sr\_1\_1?ie=UTF8\&qid=1409057181\&sr=8-1\&keywords=seraphim+rose+augustine}}, and the other by Jean Luc-Marion\footnote{\url{https://www.amazon.com/Selfs-Place-Approach-Augustine-Cultural/dp/0804762910/ref=sr\_1\_1?ie=UTF8\&qid=1409057110\&sr=8-1\&keywords=jean+luc+marion+augustine}}. Since Augustine is such a seminal watershed for both the classical and the medieval world (and therefore, the Western world), his \emph{Confessions} are an important measuring stick. Like Jean Calvin in \emph{The Institutes}, he teaches that to know yourself and to know God are in the end, sharing the same goal; unlike Calvin, he understands that they come at last to the same thing. The Self and the Self-beyond-the-Self always exist together. It might be useful to remember that the definition of a ``person\footnote{\url{http://orthosphere.org/2014/08/03/book-review-persons-the-difference-between-someone-and-something/}}'' (an individuum, not a personality) is that they have transcended themselves and recognize other Self(s). This occurs firstly by recognizing the fragmented nature of the personality and the latent existence of a higher Self, which is then a necessary goal to achieve.

Recognizing the metaphysical basis of other selves and accommodating the appearances is arguably a Western distinctive uniqueness, and it comes at a price, which is why constant contact with the Tradition has to be maintained, and explains why simply replacing the West with the East is not advisable, since the East has its own unique emphasis, also in perfect ultimate harmony with Tradition. It is nonsensical to mistake a unique emphasis and artful native form with the essence of Tradition, just as it is also nonsensical to suppose that that essence can be reached in a vacuum without developing a suitable ``vehicle'' within a given cultural Form: neither triumphalism or syncretism is sophisticated enough to do justice to the situation. Deeper Western thinkers used to do more justice to spiritual complexity.

Similar to Mouravieff's teaching that ``shocks'' or ``falls'' can actually prompt self-reflection and self-remembering, Saint Augustine states that whether he was virtuous or lost in vice, his life redounded to the praise and glory of the Lord, because the Lord's name was great, and greatly to be praised. He had discovered that his ``great question'' (who am I? a mystery to myself... ) ended in the great mystery (a God without end): the two become One.

So the basis of the \emph{Confessions} is the dual awareness that the ``I'' is an illusion, and that the ultimate resolution of this paradox would end in gnosis of the one true God. The Eastern Orthodox criticize (and to some extent rightly so) Augustine's conception of grace and his relation of it to free will, but that is not really the spiritual point of the \emph{Confessions}, something which Jean Luc-Marion's work makes perfectly clear: Augustine was teaching that the fragmented, fallen self finds its true self in the truest Self, ``that our hearts are restless till they rest in Thee''. So the inner meaning of Augustinianism is precisely in accord with what Tomberg \& Mouravieff state.


\flrightit{Posted on 2014-08-30 by Logres}

\begin{center}* * *\end{center}

\begin{footnotesize}
\begin{sffamily}

\url{https://matthewnanderson.wordpress.com/2016/03/06/ego-defense-mechanisms/}
\hfill

\end{sffamily}
\end{footnotesize}

